% Generated draft: requires mathematical and source review.
\documentclass[11pt]{article}
\usepackage[margin=25mm]{geometry}
\usepackage{fontspec}
\setmainfont{Latin Modern Roman}
\usepackage{amsmath,amssymb,mathtools,unicode-math}
\setmathfont{Latin Modern Math}
\usepackage{enumitem,needspace,xurl,hyperref}
\setlength{\parindent}{0pt}
\setlength{\parskip}{5pt}
\setlength{\emergencystretch}{3em}
\newcommand{\N}{\mathbb{N}}
\newcommand{\Z}{\mathbb{Z}}
\newcommand{\Q}{\mathbb{Q}}
\newcommand{\R}{\mathbb{R}}
\newcommand{\C}{\mathbb{C}}
\newcommand{\F}{\mathbb{F}}
\newcommand{\A}{\mathbb{A}}
\newcommand{\PP}{\mathbb{P}}
\newcommand{\E}{\mathbb{E}}
\newcommand{\eps}{\varepsilon}
\newcommand{\dd}{\,\mathrm{d}}
\begin{document}
\section*{UnsolvedMath ID: 11000178}
\textbf{Source problem code:} catalogue entry
\subsection{Definitions and mathematical statement}
Let \(\mathcal I_g\) be the Torelli group and \(\mathcal K_g\) the Johnson kernel in the source's surface convention. Fix a genus-\(g\) handlebody with boundary \(\Sigma_g\), and let \(\mathcal H_g\leq\operatorname{Mod}_g\) be the subgroup of boundary mapping classes extending over it. Put
\[A=\mathcal I_g\cap\mathcal H_g,\qquad B=\mathcal K_g\cap\mathcal H_g.\]
The double-coset sets to be studied are
\[A\backslash\mathcal I_g/A,\qquad B\backslash\mathcal K_g/B,\]
with representatives \(\varphi\) in their respective groups; \(\varphi'\) represents the same class as \(\varphi\) exactly when \(\varphi'=a\varphi b\) for suitable elements of the indicated left and right subgroup.

The problem asks for quotients of \(\mathcal I_g\) or \(\mathcal K_g\) in which the image of the appropriate handlebody intersection is tractable enough to study these sets. For a quotient \(q:\mathcal I_g\twoheadrightarrow Q\), for example, there is an induced map
\[A\backslash\mathcal I_g/A\longrightarrow q(A)\backslash Q/q(A),\qquad
A\varphi A\longmapsto q(A)q(\varphi)q(A),\]
and analogously for a quotient of \(\mathcal K_g\). A proposed construction should identify the image subgroup and what double-coset information its quotient provides. The source does not require this induced map to be injective, prescribe finite or abelian quotients, or specify a decision algorithm/completeness criterion. A quotient whose double-coset structure is easy but loses all relevant information would not identify the intended usefulness, so the adequacy criterion is an explicit gap.

The motivation concerns potential invariants of Heegaard splittings. The source contrasts this with Morita [36], described as using three-manifold invariants to understand representations of \(\mathcal I_g\) and \(\mathcal K_g\), without asking about simultaneous splitting invariants. Thus invariants of a given splitting and invariants only of its ambient three-manifold are distinct targets. No identification of the displayed double-coset sets with an unstated complete topological classification is asserted.
\subsection{Short English statement}
Find useful quotients of Torelli or Johnson kernel where the relevant handlebody-subgroup images make double cosets accessible, potentially giving Heegaard-splitting invariants. The source does not specify how much information those quotients must retain.
\subsection{Sources}
\begin{itemize}
\item Local UnsolvedMath registry, permanent ID \(11000178\), exact source code \texttt{catalogue entry}.
\item Farb, \emph{Problems on Mapping Class Groups and Related Topics}, Problem 3.7, PDF page 162: \url{https://www.math.uchicago.edu/~farb/papers/mcgbook.pdf}. Registry access 2026-07-29; pdf-font-structure extraction.
\item Morita [36] is cited for the three-manifold-invariant/representation motivation; its full entry and exact results are absent.
\end{itemize}
\end{document}
